\section{Aggregate implications}
\label{sec:aggregate_implications}

\noindent Under additional assumptions, our firm-level estimates can be aggregated to infer the total effect of leverage risk. This aggregation exercise sheds light on the macroeconomic significance of our results and helps address the following two questions: What would have happened to investment if the Supreme Court had preempted leverage risk by ruling the abrogation of gold clauses constitutional immediately after Congress passed the joint resolution in June 1933? What would have been the consequences for investment if, instead of being eliminated in early 1935, leverage risk had persisted over 1935--1936?
	
To address the first question, we abstract from general equilibrium effects and assume that, in the absence of leverage risk in the economy, all firms would invest as if they had no exposure to gold clauses (as if $d_i = 0$). Under this assumption, the rate of forgone aggregate investment due to leverage risk in 1933 and 1934 is given by
\begin{equation}
\label{eq:aggregate}
\text{Leverage risk effect}_t = \frac{\beta_1 \sum_i d_i \text{Fixed capital}_{t-1}}{\sum_i \text{Fixed capital}_{t-1}},
\end{equation}
where $i$ indexes firms, $\beta_1$ is the impact of leverage risk on investment reported in Panel A of Table \ref{tab:T4}, and $d_i$ is firm $i$'s leverage risk exposure as in (\ref{eq:d}). Similarly, the aggregate effect of the elimination of leverage risk in 1935 and 1936 can be estimated as in equation (\ref{eq:aggregate}), where $\beta_1$ is the treatment reversal estimate in Panel B of Table \ref{tab:T4}.\footnote{See the appendix for a detailed explanation of aggregate capital and investment computations.}

Table \ref{tab:T11} reports the results. The first row of Panel A shows that, in 1933 and 1934, the aggregate net investment rates of public firms were $-4.15$\% and $-2.58$\%, respectively.\footnote{Total net investment is given by $\frac{\sum_i \text{Fixed capital}_{t}}{\sum_i \text{Fixed capital}_{t-1}} - 1$.} The second row shows that our estimates of the leverage risk effect based on equation (\ref{eq:aggregate}) are $-1.28$\% and $-1.36$\% in 1933 and 1934, respectively. Thus, our estimates imply that preempting leverage risk would have cut firms' divestment by approximately one-third in 1933 and 1934. Panel B of Table \ref{tab:T11} shows that leverage risk accounts for about half of the divestment in 1933 and 1934 among $d > 0$ firms.

The last two columns of Panel A report that, in 1935 and 1936, aggregate net investment rates were -0.05\% and 1.13\%, respectively. Our estimates of the aggregate effect of the elimination of leverage risk are 1.11\% and 1.13\%, respectively. Thus, the Supreme Court's decision nearly stopped divestment in 1935 and accounts for most of the positive net investment in 1936.

Taken together, our partial equilibrium aggregation exercise suggests that the uncertainty about future leverage surrounding the abrogation of gold clauses depressed the investments of exposed public firms and therefore contributed to the slow recovery from the Great Depression. Our findings complement the literature on credit intermediation which has shown that credit constraints were quantitatively important at the beginning of the crisis \cite{Benmelech2018}.

Finally, we interpret our aggregation results as providing a lower bound for the total effect of abrogation on private investment. As previously discussed, gold clauses were also prevalent in other private contracts, including farm mortgages, public utility bonds, and railroad securities. While data limitations prevent us from directly studying the impact of abrogation on these sectors, our results suggest that the effects may have been economically significant for these firms as well.\footnote{Disaggregated data on farm mortgages with gold clauses are not available for this period. Additionally, most public utilities and railroads were not publicly traded during the 1930s, making it impossible to implement standard investment regressions for these firms.}